%%%%%%%%%%%%%%%%%%%%========== マクロ ==========%%%%%%%%%%%%%%%%%%%%
% 括弧類のマクロは，星付きで呼ぶと括弧の大きさを自動調整しない．
% 例：\paren{x}は\left(x\right)，\paren*{x}は(x)になる．

%%%%%%%%%% 用語 %%%%%%%%%%

% 用語を太字で示し，用語索引に載せる．
% 省略可能な引数には読みを書く．例：\term[しゅうごう]{集合}
\NewDocumentCommand{\term}{om}{%
	\textbf{#2}%
	\IfValueTF{#1}{\index[widx]{#1@#2}}{\index[widx]{#2}}%
}
\RenewDocumentCommand{\emph}{m}{\textbf{#1}}

%%%%%%%%%% 括弧 %%%%%%%%%%

\NewDocumentCommand{\paren}{sm}{%
	\IfBooleanTF{#1}{%
		\mathopen{} (#2) \mathclose{}%
	}{%
		\mathopen{} \left( #2 \right) \mathclose{}%
	}%
}

\NewDocumentCommand{\squarebrackets}{sm}{%
	\IfBooleanTF{#1}{%
		\mathopen{} [#2] \mathclose{}%
	}{%
		\mathopen{} \left[ #2 \right] \mathclose{}%
	}%
}

\NewDocumentCommand{\anglebrackets}{sm}{%
	\IfBooleanTF{#1}{%
		\mathopen{} \langle #2 \rangle \mathclose{}%
	}{%
		\mathopen{} \left\langle #2 \right\rangle \mathclose{}%
	}%
}

\NewDocumentCommand{\vertbracket}{sm}{%
	\IfBooleanTF{#1}{%
		\mathopen{} \lvert #2 \rvert \mathclose{}%
	}{%
		\mathopen{} \left\lvert #2 \right\rvert \mathclose{}%
	}%
}

%%%%%%%%%% 関数と集合 %%%%%%%%%%

% 関数適用．例：\apply{f}{x}はf(x)になる．
\NewDocumentCommand{\apply}{smm}{%
	\IfBooleanTF{#1}{#2\paren*{#3}}{#2\paren{#3}}%
}

% 絶対値
\NewDocumentCommand{\absolute}{sm}{%
	\IfBooleanTF{#1}{\vertbracket*{#2}}{\vertbracket{#2}}%
}

% 同値類
\NewDocumentCommand{\equivclass}{sm}{%
	\IfBooleanTF{#1}{\squarebrackets*{#2}}{\squarebrackets{#2}}%
}

% 順序対
\NewDocumentCommand{\pair}{m}{\langle #1 \rangle}

% 冪集合
\NewDocumentCommand{\powerset}{m}{\apply{\symfrak{P}}{#1}}

% 補集合
\NewDocumentCommand{\complementset}{m}{{#1}^{c}}

% 写像の制限
\NewDocumentCommand{\restriction}{}{\upharpoonright}

% 一意存在
\NewDocumentCommand{\uexists}{}{{\exists !}}

% 同型
\NewDocumentCommand{\isomorphic}{}{\cong}
